% !TEX root = main.tex
\section{Conclusions}\label{conclusions}

This simulation study demonstrates that integrating Big Five personality detection with Zurich Model-driven regulation in LLM-based conversational agents is technically feasible, achieving near-perfect implementation fidelity (98.3\% detection accuracy, 100\% regulation adherence) and a 92-percentage-point improvement in personality-specific need fulfillment (Cohen's d = 4.42, p \textless{} 0.001). The critical finding is \textbf{selective enhancement}: dramatic gains on personalization-sensitive metrics with no degradation in basic conversational quality (emotional tone: d = 0.000; relevance: d = 0.183, ns), confirming that theory-driven adaptation targets individual differences rather than generic interaction quality.

Results represent an \textbf{upper bound} under idealized conditions (boundary-condition profiles, pre-scripted dialogues, synthetic data). Real-world performance will likely be substantially lower (estimated d = 0.5--1.5). Human subject validation with moderate personality profiles, longitudinal tracking, user-reported outcomes, and bias audits is mandatory before real-world deployment. The detection $\rightarrow$ regulation $\rightarrow$ evaluation framework serves as an augmentation layer for existing digital mental health interventions, adaptable to alternative personality theories and support domains, providing a methodological template for AI-augmented personalized mental health support grounded in psychological theory.

